\FigureLayoutDeclare{img/2007/jiangxi_liberal/q11_fig1.png}{14.0cm}{31bp 26bp 30bp 23bp}
\FigureLayoutDeclare{img/2007/jiangxi_liberal/q16_fig1.png}{7.0cm}{22bp 26bp 64bp 3bp}
\FigureLayoutDeclare{img/2007/jiangxi_liberal/q18_fig1.png}{10.0cm}{10bp 9bp 9bp 10bp}

\chapter{2007年江西卷（文）}

\section{单选题}

\begin{problem}
若集合 \(M = \{0, 1\}\)，\(I = \{0, 1, 2, 3, 4, 5\}\)，则 \(\complement_I M\) 为（\quad）

\choices
    {\(\{0, 1\}\)}
    {\(\{2, 3, 4, 5\}\)}
    {\(\{0, 2, 3, 4, 5\}\)}
    {\(\{1, 2, 3, 4, 5\}\)}

\begin{answer}
B
\end{answer}

\begin{solution}
由补集的定义，\(\complement_I M\) 由全集 \(I\) 中不属于 \(M\) 的元素组成，因此
\(\complement_I M=\{2,3,4,5\}\). 故选 B.
\end{solution}
\end{problem}

\begin{problem}
函数 \( y = 5\tan (2x + 1) \) 的最小正周期为（\quad）

\choices
    {\(\frac{\pi}{4}\)}
    {\(\frac{\pi}{2}\)}
    {\(\pi\)}
    {\(2\pi\)}

\begin{answer}
B
\end{answer}

\begin{solution}
正切函数的最小正周期为 \(\pi\). 对于 \(y=5\tan(2x+1)\)，令自变量增加 \(T\)，要使函数值重复，需满足 \(2T=\pi\). 因而最小正周期为 \(T=\frac{\pi}{2}\)，故选 B.
\end{solution}
\end{problem}

\begin{problem}
函数 \( f(x) = \lg \frac{1-x}{x-4} \) 的定义域为（\quad）

\choices
    {\((1, 4)\)}
    {\([1,4)\)}
    {\((-\infty, 1) \cup (4, +\infty)\)}
    {\((-\infty, 1] \cup (4, +\infty)\)}

\begin{answer}
A
\end{answer}

\begin{solution}
对数的真数必须为正，即
\(\frac{1-x}{x-4}>0\). 临界点为 \(x=1\) 和 \(x=4\). 当 \(x<1\) 时，分子为正、分母为负；当 \(1<x<4\) 时，分子和分母均为负；当 \(x>4\) 时，分子为负、分母为正. 因此真数为正的范围是 \(1<x<4\)，定义域为 \((1,4)\)，故选 A.
\end{solution}
\end{problem}

\begin{problem}
若 \(\tan \alpha = 3\)，\(\tan \beta = \frac{4}{3}\)，则 \(\tan (\alpha - \beta)\) 等于（\quad）

\choices
    {\(-3\)}
    {\( -\frac{1}{3} \)}
    {\(3\)}
    {\( \frac{1}{3} \)}

\begin{answer}
D
\end{answer}

\begin{solution}
由正切的差角公式，
\[
\tan(\alpha-\beta)=\frac{\tan\alpha-\tan\beta}{1+\tan\alpha\tan\beta}=\frac{3-\frac{4}{3}}{1+3\cdot\frac{4}{3}}=\frac{\frac{5}{3}}{5}=\frac{1}{3}.
\]
故选 D.
\end{solution}
\end{problem}

\begin{problem}
设 \((x^{2} + 1)(2x + 1)^{9} = a_{0} + a_{1}(x + 2) + a_{2}(x + 2)^{2} + \dots +a_{11}(x + 2)^{11}\)，则 \(a_0 + a_1 + a_2 + \dots +a_{11}\) 的值为（\quad）

\choices
    {\(-2\)}
    {\(-1\)}
    {\(1\)}
    {\(2\)}

\begin{answer}
A
\end{answer}

\begin{solution}
令
\(P(x)=(x^2+1)(2x+1)^9\). 题设给出的展开式是在 \(x+2\) 的幂次下展开的，令 \(x+2=1\)，即取 \(x=-1\)，得到
\[
a_0+a_1+\cdots+a_{11}=P(-1)=\bigl((-1)^2+1\bigr)(-2+1)^9=2\cdot(-1)=-2.
\]
故选 A.
\end{solution}
\end{problem}

\begin{problem}
一袋中装有大小相同，编号分别为 \(1\)，\(2\)，\(3\)，\(4\)，\(5\)，\(6\)，\(7\)，\(8\) 的八个球，从中有放回地每次取一个球，共取 \(2\) 次，则取得两个球的编号和不小于 \(15\) 的概率为（\quad）

\choices
    {\(\frac{1}{32}\)}
    {\(\frac{1}{64}\)}
    {\(\frac{3}{32}\)}
    {\(\frac{3}{64}\)}

\begin{answer}
D
\end{answer}

\begin{solution}
由于每次取球后放回，两个球的编号组成的有序对共有 \(8\times8=64\) 个，且等可能. 和不小于 \(15\) 的有序对只有
\((7,8)\)、\((8,7)\) 和 \((8,8)\)，共 \(3\) 个. 所求概率为 \(\frac{3}{64}\)，故选 D.
\end{solution}
\end{problem}

\begin{problem}
连接抛物线 \(x^{2} = 4y\) 的焦点 \(F\) 与点 \(M(1,0)\) 所得的线段与抛物线交于点 \(A\). 设点 \(O\) 为坐标原点，则三角形 \(OAM\) 的面积为（\quad）

\choices
    {\(-1 + \sqrt{2}\)}
    {\(\frac{3}{2} -\sqrt{2}\)}
    {\(1 + \sqrt{2}\)}
    {\(\frac{3}{2} +\sqrt{2}\)}

\begin{answer}
B
\end{answer}

\begin{solution}
抛物线 \(x^2=4y\) 的焦点为 \(F(0,1)\). 直线 \(FM\) 经过 \(F(0,1)\) 和 \(M(1,0)\)，方程为 \(y=1-x\). 联立抛物线方程，得
\[
x^2=4(1-x)\Longleftrightarrow x^2+4x-4=0,
\]
所以交点 \(A\) 在线段 \(FM\) 上，对应的横坐标为 \(x_A=-2+2\sqrt{2}\)，纵坐标为
\(y_A=1-x_A=3-2\sqrt{2}\). 以 \(OM\) 为底，\(OM=1\)，三角形 \(OAM\) 的高为 \(y_A\)，故
\[
S_{\triangle OAM}=\frac{1}{2}\cdot1\cdot(3-2\sqrt{2})=\frac{3}{2}-\sqrt{2}.
\]
故选 B.
\end{solution}
\end{problem}

\begin{problem}
若 \(0 < x < \frac{\pi}{2}\)，则下列命题正确的是（\quad）

\choices
    {\(\sin x < \frac{2}{\pi}x\)}
    {\(\sin x > \frac{2}{\pi}x\)}
    {\(\sin x < \frac{3}{\pi}x\)}
    {\(\sin x > \frac{3}{\pi}x\)}

\begin{answer}
B
\end{answer}

\begin{solution}
函数 \(y=\sin x\) 在 \(\left[0,\frac{\pi}{2}\right]\) 上满足
\(y''=-\sin x<0\)，所以图象严格向下凹. 连接端点 \(\left(0,0\right)\) 和 \(\left(\frac{\pi}{2},1\right)\) 的弦所在直线为
\(y=\frac{2}{\pi}x\). 严格凹函数的图象在弦的上方，因此当 \(0<x<\frac{\pi}{2}\) 时，
\[
\sin x>\frac{2}{\pi}x.
\]
故选 B.
\end{solution}
\end{problem}

\begin{problem}
四面体 \(ABCD\) 外接球球心在 \(CD\) 上，且 \(CD = 2\)，\(AB = \sqrt{3}\)，在外接球面上两点 \(A\)，\(B\) 间的球面距离是（\quad）

\choices
    {\(\frac{\pi}{6}\)}
    {\(\frac{\pi}{3}\)}
    {\(\frac{2\pi}{3}\)}
    {\(\frac{5\pi}{6}\)}

\begin{answer}
C
\end{answer}

\begin{solution}
外接球球心在弦 \(CD\) 所在直线上，且 \(C,D\) 都在球面上，所以 \(CD\) 是球的直径. 由 \(CD=2\)，得外接球半径 \(R=1\).

设球心为 \(O\)，且小于或等于 \(\pi\) 的圆心角 \(\angle AOB=\varphi\). 弦长公式给出
\[
AB=2R\sin\frac{\varphi}{2}=\sqrt{3},
\]
从而 \(\sin\frac{\varphi}{2}=\frac{\sqrt{3}}{2}\)，所以 \(\varphi=\frac{2\pi}{3}\). 球面距离等于对应大圆弧长，故
\[
R\varphi=1\cdot\frac{2\pi}{3}=\frac{2\pi}{3}.
\]
故选 C.
\end{solution}
\end{problem}

\begin{problem}
设 \( p: f(x) = x^3 + 2x^2 + mx + 1 \) 在 \((- \infty, +\infty)\) 内单调递增，\( q: m \geqslant \frac{4}{3} \)，则 \( p \) 是 \( q \) 的（\quad）

\choices
    {充分不必要条件}
    {必要不充分条件}
    {充分必要条件}
    {既不充分也不必要条件}

\begin{answer}
C
\end{answer}

\begin{solution}
函数 \(f(x)=x^3+2x^2+mx+1\) 的导数为
\[
f'(x)=3x^2+4x+m=3\left(x+\frac{2}{3}\right)^2+m-\frac{4}{3}.
\]
要使 \(f(x)\) 在整个实数轴上单调递增，必须且只需 \(f'(x)\geqslant0\) 对一切 \(x\in\R\) 成立，即
\(m-\frac{4}{3}\geqslant0\)，也就是 \(m\geqslant\frac{4}{3}\). 因此命题 \(p\) 与命题 \(q\) 等价，\(p\) 是 \(q\) 的充分必要条件，故选 C.
\end{solution}
\end{problem}

\begin{problem}
四位好朋友在一次聚会上，他们按照各自的爱好选择了形状不同，内空高度相等，杯口半径相等的圆口酒杯，如图所示，盛满酒后他们约定：先各自饮杯中酒的一半. 设剩余酒的高度从左到右依次为 \(h_1\)，\(h_2\)，\(h_3\)，\(h_4\)，则它们的大小关系正确的是（\quad）

\bitmapfigure[max width=0.84\linewidth]{img/2007/jiangxi_liberal/q11_fig1.png}

\choices
    {\(h_2 > h_1 > h_4\)}
    {\(h_1 > h_2 > h_3\)}
    {\(h_3 > h_2 > h_4\)}
    {\(h_2 > h_4 > h_1\)}

\begin{answer}
A
\end{answer}

\begin{solution}
设四个杯子的内空高度均为 \(H\)，以杯底为高度 \(0\)，记高度为 \(t\) 处的水平截面积为 \(S(t)\). 饮去一半后，剩余酒的体积等于该杯总容积的一半.

第四个杯子是圆柱，截面积处处相等，所以剩余高度正好为
\(h_4=\frac{H}{2}\). 第一个杯子和第二个杯子的截面积都随高度增加而增加，因此上半部单位高度的容积大于下半部，达到总容积一半时高度均超过 \(\frac{H}{2}\)，特别有 \(h_1>h_4\).

比较前两个杯子的轮廓，第二个杯子在下部更细，容积相对更多地集中在上部；第一个杯子较早变宽，容积相对更多地分布在下部. 因而取得各自总容积的一半时，第二个杯子需要达到更高的高度，即 \(h_2>h_1\). 所以
\[
h_2>h_1>h_4,
\]
故选 A.
\end{solution}
\end{problem}

\begin{problem}
设椭圆 \(\frac{x^2}{a^2} + \frac{y^2}{b^2} = 1\)（\(a > b > 0\)）的离心率为 \(e = \frac{1}{2}\)，右焦点为 \(F(c,0)\). 方程 \(ax^2 + bx - c = 0\) 的两个实根分别为 \(x_1\) 和 \(x_2\)，则点 \(P(x_1, x_2)\)（\quad）

\choices
    {必在圆 \(x^{2} + y^{2} = 2\) 上}
    {必在圆 \(x^{2} + y^{2} = 2\) 外}
    {必在圆 \(x^{2} + y^{2} = 2\) 内}
    {以上三种情形都有可能}

\begin{answer}
C
\end{answer}

\begin{solution}
椭圆的离心率满足 \(e=\frac{c}{a}=\frac{1}{2}\)，所以 \(c=\frac{a}{2}\)，并且
\(b^2=a^2-c^2=\frac{3a^2}{4}\)，从而 \(b=\frac{\sqrt{3}}{2}a\).

由方程 \(ax^2+bx-c=0\) 的韦达定理，
\[
x_1+x_2=-\frac{b}{a}=-\frac{\sqrt{3}}{2},
\qquad
x_1x_2=-\frac{c}{a}=-\frac{1}{2}.
\]
因此点 \(P(x_1,x_2)\) 到原点的距离平方为
\[
x_1^2+x_2^2=(x_1+x_2)^2-2x_1x_2=\frac{3}{4}+1=\frac{7}{4}<2.
\]
所以点 \(P\) 必在圆 \(x^2+y^2=2\) 内，故选 C.
\end{solution}
\end{problem}

\section{填空题}

\begin{problem}
在平面直角坐标系中，正方形 \(OABC\) 的对角线 \(OB\) 的两端点分别为 \(O(0,0)\)，\(B(1,1)\)，则 \(\overrightarrow{AB} \cdot \overrightarrow{AC} =\)\fillinblank{}.

\begin{answer}
\(1\)
\end{answer}

\begin{solution}
不妨取正方形的另外两个顶点为 \(A(1,0)\)、\(C(0,1)\)；交换 \(A,C\) 的位置不会改变所求内积. 则
\[
\overrightarrow{AB}=B-A=(0,1),
\qquad
\overrightarrow{AC}=C-A=(-1,1).
\]
所以
\[
\overrightarrow{AB}\cdot\overrightarrow{AC}=0\cdot(-1)+1\cdot1=1.
\]
\end{solution}
\end{problem}

\begin{problem}
已知等差数列 \(\{a_{n}\}\) 的前 \(n\) 项和为 \(S_{n}\)，若 \(S_{12} = 21\)，则 \(a_{2} + a_{5} + a_{8} + a_{11} =\)\fillinblank{}.

\begin{answer}
\(7\)
\end{answer}

\begin{solution}
设等差数列的首项为 \(a_1\)，公差为 \(d\)，则 \(a_k=a_1+(k-1)d\). 由等差数列前 \(12\) 项和公式，
\[
S_{12}=\frac{12(a_1+a_{12})}{2}=6(a_1+a_{12})=21,
\]
所以
\[
a_1+a_{12}=2a_1+11d=\frac{7}{2}.
\]
另一方面，
\[
a_2+a_{11}=2a_1+11d=\frac{7}{2},
\qquad
a_5+a_8=2a_1+11d=\frac{7}{2}.
\]
因此
\[
a_2+a_5+a_8+a_{11}=\frac{7}{2}+\frac{7}{2}=7.
\]
\end{solution}
\end{problem}

\begin{problem}
已知函数 \( y = f(x) \) 存在反函数 \( y = f^{-1}(x) \)，若函数 \( y = f(1 + x) \) 的图象经过点 \((3,1)\)，则函数 \( y = f^{-1}(x) \) 的图象必经过点\fillinblank{}.

\begin{answer}
\((1,4)\)
\end{answer}

\begin{solution}
函数 \(y=f(1+x)\) 的图象经过 \((3,1)\)，代入得
\[
f(1+3)=f(4)=1.
\]
反函数图象由原函数图象关于直线 \(y=x\) 对称，原图象上的点 \((4,1)\) 对称后变为 \((1,4)\). 因此函数 \(y=f^{-1}(x)\) 的图象必经过 \((1,4)\).
\end{solution}
\end{problem}

\begin{problem}
如图，正方体 \(ABCD-A_{1}B_{1}C_{1}D_{1}\) 的棱长为 \(1\)，过点 \(A\) 作平面 \(A_{1}BD\) 的垂线，垂足为点 \(H\). 有下列四个命题：

\begin{circlelist}
\item 点 \(H\) 是 \(\triangle A_{1}BD\) 的垂心；
\item \(AH \perp\) 平面 \(CB_{1}D_{1}\)；
\item 二面角 \(C - B_{1}D_{1} - C_{1}\) 的正切值为 \(\sqrt{2}\)；
\item 点 \(H\) 到平面 \(A_{1}B_{1}C_{1}D_{1}\) 的距离为 \(\frac{3}{4}\).
\end{circlelist}

其中真命题的代号是\fillinblank{}.（写出所有真命题的代号）

\bitmapfigure[max width=0.55\linewidth]{img/2007/jiangxi_liberal/q16_fig1.png}

\begin{answer}
A、B、C
\end{answer}

\begin{solution}
建立空间直角坐标系，取
\[
A_1(0,0,0),\quad B_1(1,0,0),\quad C_1(1,1,0),\quad D_1(0,1,0),
\]
并令
\[
A(0,0,1),\quad B(1,0,1),\quad C(1,1,1),\quad D(0,1,1).
\]
平面 \(A_1BD\) 的方程为 \(z=x+y\)，其一个法向量为 \((-1,-1,1)\). 从点 \(A\) 向该平面作垂线，垂足为
\[
H=A-\frac{1}{(-1)^2+(-1)^2+1^2}(-1,-1,1)=\left(\frac{1}{3},\frac{1}{3},\frac{2}{3}\right).
\]

先验证命题 A. 三角形 \(A_1BD\) 的三条边方向向量可取 \(\overrightarrow{BD}=(-1,1,0)\)、\(\overrightarrow{A_1D}=(0,1,1)\) 和 \(\overrightarrow{A_1B}=(1,0,1)\). 有
\[
\overrightarrow{A_1H}=\left(\frac{1}{3},\frac{1}{3},\frac{2}{3}\right),
\quad
\overrightarrow{BH}=\left(-\frac{2}{3},\frac{1}{3},-\frac{1}{3}\right),
\quad
\overrightarrow{DH}=\left(\frac{1}{3},-\frac{2}{3},-\frac{1}{3}\right).
\]
于是
\[
\overrightarrow{A_1H}\cdot\overrightarrow{BD}=0,
\quad
\overrightarrow{BH}\cdot\overrightarrow{A_1D}=0,
\quad
\overrightarrow{DH}\cdot\overrightarrow{A_1B}=0.
\]
故 \(H\) 同时在三条高上，是 \(\triangle A_1BD\) 的垂心，命题 A 正确.

再验证命题 B. 平面 \(CB_1D_1\) 内的两个不共线方向向量可取 \(\overrightarrow{B_1D_1}=(-1,1,0)\) 和 \(\overrightarrow{B_1C}=(0,1,1)\)，其法向量可取 \((1,1,-1)\). 而
\[
\overrightarrow{AH}=\left(\frac{1}{3},\frac{1}{3},-\frac{1}{3}\right)
\]
与该法向量平行，所以 \(AH\perp\) 平面 \(CB_1D_1\)，命题 B 正确.

验证命题 C. 二面角 \(C-B_1D_1-C_1\) 是平面 \(CB_1D_1\) 与平面 \(C_1B_1D_1\) 所成的角. 两个平面的法向量可分别取 \(\bs n_1=(1,1,-1)\) 和 \(\bs n_2=(0,0,1)\). 设二面角的锐角为 \(\varphi\)，则
\[
\cos\varphi=\frac{|\bs n_1\cdot\bs n_2|}{|\bs n_1||\bs n_2|}=\frac{1}{\sqrt{3}},
\]
从而
\[
\tan\varphi=\frac{\sqrt{1-\cos^2\varphi}}{\cos\varphi}=\frac{\sqrt{\frac{2}{3}}}{\frac{1}{\sqrt{3}}}=\sqrt{2}.
\]
故命题 C 正确.

最后，平面 \(A_1B_1C_1D_1\) 的方程为 \(z=0\)，所以点 \(H\) 到该平面的距离为
\[
\frac{2}{3}\ne\frac{3}{4},
\]
所以命题 D 错误，前三个命题为真，结论与填空处一致.
\end{solution}
\end{problem}

\section{解答题}

\begin{problem}
已知函数 \( f(x)=\begin{cases}
cx+1,&0<x<c,\\
2^{-\frac{x}{c^{2}}}+1,&c\leqslant x<1,
\end{cases} \) 在区间 \( (0,1) \) 内连续，且 \( f(c^{2})=\frac{9}{8}. \)

\begin{enumerate}
\item 求常数 \(c\) 的值；
\item 解不等式 \( f(x) > \frac{\sqrt{2}}{8} + 1 \).
\end{enumerate}

\begin{solution}
\begin{enumerate}
\item 由 \(f(c^2)\) 有意义且定义域为 \((0,1)\)，得 \(c^2\in(0,1)\). 若 \(-1<c<0\)，则 \(c<c^2<1\)，由第二段可得 \(f(c^2)=2^{-c^2/c^2}+1=\frac{3}{2}\ne\frac{9}{8}\)，所以 \(0<c<1\). 因此 \(c^2<c\)，从而 \(f(c^2)=c\cdot c^2+1=c^3+1=\frac{9}{8}\)，解得 \(c=\frac{1}{2}\). 此时在分界点 \(x=c\) 处，左侧极限为 \(c^2+1=\frac{5}{4}\)，右侧函数值为 \(2^{-1/c}+1=\frac{5}{4}\)，满足连续性条件.

\item 当 \(c=\frac{1}{2}\) 时，不等式为 \(f(x)>1+\frac{\sqrt{2}}{8}\). 当 \(0<x<\frac{1}{2}\) 时，\(f(x)=\frac{x}{2}+1\)，所以 \(\frac{x}{2}+1>1+\frac{\sqrt{2}}{8}\Longleftrightarrow x>\frac{\sqrt{2}}{4}\)，得到 \(\frac{\sqrt{2}}{4}<x<\frac{1}{2}\). 当 \(\frac{1}{2}\leqslant x<1\) 时，\(f(x)=2^{-4x}+1\)，且 \(\frac{\sqrt{2}}{8}=2^{-5/2}\)，所以 \(2^{-4x}>2^{-5/2}\Longleftrightarrow x<\frac{5}{8}\)，得到 \(\frac{1}{2}\leqslant x<\frac{5}{8}\). 合并两段，原不等式的解集为 \(\left(\frac{\sqrt{2}}{4},\frac{5}{8}\right)\).
\end{enumerate}
\end{solution}
\end{problem}

\begin{problem}
如图，函数 \( y = 2\cos (\omega x + \theta) \)（\(x \in \R\)，\(0 \leqslant \theta \leqslant \frac{\pi}{2}\)）的图象与 \(y\) 轴交于点 \( (0, \sqrt{3}) \)，且在该点处切线的斜率为 \(-2\).

\begin{enumerate}
\item 求 \( \theta \) 和 \( \omega \) 的值；
\item 已知点 \(A\left(\frac{\pi}{2}, 0\right)\)，点 \(P\) 是该函数图象上的一点，点 \(Q(x_0, y_0)\) 是 \(PA\) 的中点，当 \(y_0 = \frac{\sqrt{3}}{2}\)，\(x_0 \in \left[\frac{\pi}{2}, \pi\right]\) 时，求 \(x_0\) 的值.
\end{enumerate}

\bitmapfigure[max width=0.55\linewidth]{img/2007/jiangxi_liberal/q18_fig1.png}

\begin{solution}
\begin{enumerate}
\item 由图象经过 \((0,\sqrt{3})\)，得 \(2\cos\theta=\sqrt{3}\). 由于 \(0\leqslant\theta\leqslant\frac{\pi}{2}\)，所以 \(\theta=\frac{\pi}{6}\). 函数的导数为 \(y'=-2\omega\sin(\omega x+\theta)\)，由 \(y\) 轴交点处的切线斜率为 \(-2\)，得 \(-2\omega\sin\theta=-2\)，即 \(\omega\sin\frac{\pi}{6}=1\)，故 \(\omega=2\).

\item 设 \(P(x_1,y_1)\). 因为 \(Q\) 是 \(PA\) 的中点，所以 \(x_0=\frac{x_1+\frac{\pi}{2}}{2}\)，\(y_0=\frac{y_1}{2}\). 由 \(y_0=\frac{\sqrt{3}}{2}\) 得 \(y_1=\sqrt{3}\)，再由 \(P\) 在函数图象上得 \(2\cos\left(2x_1+\frac{\pi}{6}\right)=\sqrt{3}\)，即 \(\cos\left(2x_1+\frac{\pi}{6}\right)=\frac{\sqrt{3}}{2}\). 由 \(x_0\in\left[\frac{\pi}{2},\pi\right]\) 及 \(x_1=2x_0-\frac{\pi}{2}\)，知 \(x_1\in\left[\frac{\pi}{2},\frac{3\pi}{2}\right]\)，于是 \(2x_1+\frac{\pi}{6}\in\left[\frac{7\pi}{6},\frac{19\pi}{6}\right]\). 在此区间内满足余弦值为 \(\frac{\sqrt{3}}{2}\) 的角只有 \(\frac{11\pi}{6}\) 和 \(\frac{13\pi}{6}\)，因此 \(x_1=\frac{5\pi}{6}\) 或 \(x_1=\pi\). 代回中点公式，分别得到 \(x_0=\frac{2\pi}{3}\) 或 \(x_0=\frac{3\pi}{4}\).
\end{enumerate}
\end{solution}
\end{problem}

\begin{problem}
栽培甲，乙两种果树，先要培育成苗，然后再进行移栽. 已知甲，乙两种果树成苗的概率分别为 \(0.6\)，\(0.5\)，移栽后成活的概率分别为 \(0.7\)，\(0.9\).

\begin{enumerate}
\item 求甲，乙两种果树至少有一种果树成苗的概率；
\item 求恰好一种果树能栽培成苗且移栽成活的概率.
\end{enumerate}

\begin{solution}
\begin{enumerate}
\item 设事件 \(A\)，\(B\) 分别表示甲、乙果树成苗. 按题意各阶段的成败相互独立，所以 \(P(A\cup B)=1-P(\overline{A}\cap\overline{B})=1-(1-0.6)(1-0.5)=0.8\).
\item 设 \(C\)，\(D\) 分别表示甲、乙果树成苗且移栽成活，则 \(P(C)=0.6\times0.7=0.42\)，\(P(D)=0.5\times0.9=0.45\). 恰好一种果树成苗且移栽成活的概率为 \(P(C\cap\overline{D})+P(\overline{C}\cap D)=0.42(1-0.45)+(1-0.42)0.45=0.492\).
\end{enumerate}
\end{solution}
\end{problem}

\begin{problem}
如图是一个直三棱柱（以 \(A_{1}B_{1}C_{1}\) 为底面）被一平面所截得到的几何体. 截面为 \(ABC\). 已知 \(A_{1}B_{1} = B_{1}C_{1} = 1\)，\(\angle A_{1}B_{1}C_{1} = 90^{\circ}\)，\(AA_{1} = 4\)，\(BB_{1} = 2\)，\(CC_{1} = 3\).

\begin{enumerate}
\item 设点 \(O\) 是 \(AB\) 的中点，证明：\(OC \parallel\) 平面 \(A_{1}B_{1}C_{1}\)；
\item 求 \(AB\) 与平面 \(AA_{1}C_{1}C\) 所成的角的大小；
\item 求此几何体的体积.
\end{enumerate}

\bitmapfigure[max width=0.52\linewidth]{img/2007/jiangxi_liberal/q20_fig1.png}

\begin{solution}
\begin{enumerate}
\item 建立空间直角坐标系，取 \(A_1(0,0,0)\)，\(B_1(1,0,0)\)，\(C_1(1,1,0)\)，并令竖直方向为 \(z\) 轴正方向，则 \(A(0,0,4)\)，\(B(1,0,2)\)，\(C(1,1,3)\). \(O\) 是 \(AB\) 的中点，所以 \(O\left(\frac{1}{2},0,3\right)\)，从而 \(\overrightarrow{OC}=\left(\frac{1}{2},1,0\right)\). 平面 \(A_1B_1C_1\) 的方程为 \(z=0\)，而 \(\overrightarrow{OC}\) 的竖直分量为 \(0\)，故 \(OC\parallel\) 平面 \(A_1B_1C_1\).

\item 平面 \(AA_1C_1C\) 上的四个点都满足 \(x=y\)，故该平面的一个法向量为 \(\bs{n}=(1,-1,0)\). 取 \(\overrightarrow{AB}=(1,0,-2)\)，设 \(AB\) 与平面 \(AA_1C_1C\) 所成的角为 \(\varphi\)，则 \(\sin\varphi=\frac{|\overrightarrow{AB}\cdot\bs{n}|}{|\overrightarrow{AB}|\,|\bs{n}|}=\frac{1}{\sqrt{5}\sqrt{2}}=\frac{1}{\sqrt{10}}\). 因此 \(\cos\varphi=\frac{3}{\sqrt{10}}\)，从而 \(\tan\varphi=\frac{1}{3}\)，即 \(\varphi=\arctan\frac{1}{3}\).

\item 该几何体可分成三个互不重叠的四面体 \(A_1B_1C_1A\)，\(B_1C_1AB\) 和 \(C_1ABC\). 第一个四面体的底面 \(A_1B_1C_1\) 面积为 \(\frac{1}{2}\)，高为 \(4\)，所以 \(V_{A_1B_1C_1A}=\frac{1}{3}\cdot\frac{1}{2}\cdot4=\frac{2}{3}\). 对于四面体 \(B_1C_1AB\)，取 \(B_1\) 为基点，有 \(\overrightarrow{B_1C_1}=(0,1,0)\)，\(\overrightarrow{B_1A}=(-1,0,4)\)，\(\overrightarrow{B_1B}=(0,0,2)\)，所以

\[
V_{B_1C_1AB}=\frac{1}{6}\left|\bigl(\overrightarrow{B_1C_1}\times\overrightarrow{B_1A}\bigr)\cdot\overrightarrow{B_1B}\right|=\frac{1}{3}.
\]

对于四面体 \(C_1ABC\)，取 \(C_1\) 为基点，有 \(\overrightarrow{C_1A}=(-1,-1,4)\)，\(\overrightarrow{C_1B}=(0,-1,2)\)，\(\overrightarrow{C_1C}=(0,0,3)\)，所以

\[
V_{C_1ABC}=\frac{1}{6}\left|\bigl(\overrightarrow{C_1A}\times\overrightarrow{C_1B}\bigr)\cdot\overrightarrow{C_1C}\right|=\frac{1}{2}.
\]

故几何体的体积为 \(V=\frac{2}{3}+\frac{1}{3}+\frac{1}{2}=\frac{3}{2}\).
\end{enumerate}
\end{solution}
\end{problem}

\begin{problem}
设 \(\{a_{n}\}\) 为等比数列，\(a_1 = 1\)，\(a_2 = 3\).

\begin{enumerate}
\item 求最小的自然数 \(n\)，使 \(a_n \geqslant 2007\)；
\item 求和：\(T_{2n} = \frac{1}{a_1} - \frac{2}{a_2} + \frac{3}{a_3} - \dots - \frac{2n}{a_{2n}}\).
\end{enumerate}

\begin{solution}
\begin{enumerate}
\item 等比数列的公比为 \(q=\frac{a_2}{a_1}=3\)，所以 \(a_n=3^{n-1}\). 由于 \(3^6=729<2007\leqslant2187=3^7\)，故 \(n-1\geqslant7\)，满足条件的最小自然数为 \(n=8\).

\item 由 \(a_k=3^{k-1}\)，有 \(T_{2n}=\sum_{k=1}^{2n}k\left(-\frac{1}{3}\right)^{k-1}\). 令 \(r\in\R\)，由有限等比数列求和公式

\[
1+r+r^2+\dots+r^{2n}=\frac{1-r^{2n+1}}{1-r}.
\]

两边对 \(r\) 求导，得

\[
1+2r+3r^2+\dots+2nr^{2n-1}=\frac{1-(2n+1)r^{2n}+2nr^{2n+1}}{(1-r)^2}.
\]

令 \(r=-\frac{1}{3}\)，则

\[
T_{2n}=\frac{1-\frac{2n+1}{3^{2n}}-\frac{2n}{3^{2n+1}}}{\left(1+\frac{1}{3}\right)^2}=\frac{9}{16}\left(1-\frac{8n+3}{3^{2n+1}}\right)=\frac{9}{16}-\frac{8n+3}{16\cdot3^{2n-1}}.
\]
\end{enumerate}
\end{solution}
\end{problem}

\begin{problem}
设动点 \(P\) 到两定点 \(F_{1}(-1,0)\) 和 \(F_{2}(1,0)\) 的距离分别为 \(d_{1}\) 和 \(d_{2}\)，\(\angle F_{1}PF_{2} = 2\theta\)，且存在常数 \(\lambda\)（\(0 < \lambda < 1\)），使得 \(d_{1}d_{2}\sin^{2}\theta = \lambda\).

\begin{enumerate}
\item 证明：动点 \(P\) 的轨迹 \(C\) 为双曲线，并求出 \(C\) 的方程；
\item 如图，过点 \(F_{2}\) 的直线与双曲线 \(C\) 的右支交于 \(A\)，\(B\) 两点. 问：是否存在 \(\lambda\)，使 \(\triangle F_{1}AB\) 是以点 \(B\) 为直角顶点的等腰直角三角形？若存在，求出 \(\lambda\) 的值；若不存在，说明理由.
\end{enumerate}

\bitmapfigure[max width=0.55\linewidth]{img/2007/jiangxi_liberal/q22_fig1.png}

\begin{solution}
\begin{enumerate}
\item 设 \(P(x,y)\). 在三角形 \(F_1PF_2\) 中，\(F_1F_2=2\)，且 \(\angle F_1PF_2=2\theta\)，由余弦定理得

\[
4=d_1^2+d_2^2-2d_1d_2\cos 2\theta.
\]

因此

\[
(d_1-d_2)^2=d_1^2+d_2^2-2d_1d_2=4-2d_1d_2(1-\cos2\theta)=4-4d_1d_2\sin^2\theta=4(1-\lambda).
\]

于是 \(|d_1-d_2|=2\sqrt{1-\lambda}\). 令 \(a=\sqrt{1-\lambda}\)，则 \(0<a<1=\frac{F_1F_2}{2}\)，所以 \(P\) 的轨迹是以 \(F_1\)，\(F_2\) 为焦点、实轴长为 \(2a\) 的双曲线. 由于焦距的一半为 \(c=1\)，虚半轴的平方为 \(b^2=c^2-a^2=1-(1-\lambda)=\lambda\)，故轨迹 \(C\) 的方程为

\[
\frac{x^2}{1-\lambda}-\frac{y^2}{\lambda}=1.
\]

在右支上 \(d_1-d_2=2a\)，又因 \(d_1^2-d_2^2=4x\)，所以 \(d_1+d_2=\frac{2x}{a}\)，即 \(d_2=\frac{x}{a}-a\). 代入 \(d_2^2=(x-1)^2+y^2\)，可得

\[
\left(\frac{x}{a}-a\right)^2=(x-1)^2+y^2
\quad\Longleftrightarrow\quad
\frac{x^2}{a^2}-\frac{y^2}{1-a^2}=1.
\]

这与上式一致，故方程确实给出了该轨迹.

\item 设 \(h=BF_1\)，\(r=BF_2\)，\(s=AF_2\). 直线 \(AB\) 过焦点 \(F_2\)，并与双曲线的右支交于 \(A\)，\(B\) 两点，所以 \(F_2\) 在 \(AB\) 之间. 由双曲线定义，\(h-r=2a\)，\(AF_1-s=2a\). 若 \(\triangle F_1AB\) 是以 \(B\) 为直角顶点的等腰直角三角形，则 \(AB=BF_1=h\)，且 \(h=s+r\). 结合 \(h-r=2a\)，得 \(s=2a\). 又因为 \(AF_1=\sqrt{2}h\)，所以 \(\sqrt{2}h-s=2a\Longleftrightarrow\sqrt{2}h-2a=2a\). 从而 \(h=2\sqrt{2}a\)，\(r=h-2a=2a(\sqrt{2}-1)\).

由于 \(\angle F_1BF_2=90^{\circ}\)，三角形 \(F_1BF_2\) 的斜边为 \(F_1F_2=2\)，所以 \(h^2+r^2=4\)，即 \(4a^2(5-2\sqrt{2})=4\). 因此 \(a^2=\frac{1}{5-2\sqrt{2}}\)，从而

\[
\lambda=1-a^2=\frac{4-2\sqrt{2}}{5-2\sqrt{2}}=\frac{12-2\sqrt{2}}{17}.
\]

该值确实可实现. 对此 \(\lambda\)，取 \(B\left(\sqrt{1-\lambda^2},\lambda\right)\). 因为 \(x_B^2+y_B^2=1\)，且 \(\frac{x_B^2}{1-\lambda}-\frac{y_B^2}{\lambda}=\frac{1-\lambda^2}{1-\lambda}-\lambda=1\)，所以 \(B\) 在双曲线 \(C\) 的右支上，并且 \(BF_1\perp BF_2\). 对上面求得的 \(a\)，由前面的结果有 \(h=2\sqrt{2}a\)，\(r=2a(\sqrt{2}-1)\). 在 \(F_2B\) 的反向延长线上取点 \(A\)，使 \(AF_2=2a\). 则 \(AB=AF_2+BF_2=2a+r=h=BF_1\)，且 \(AB\perp BF_1\). 又 \(AF_1=\sqrt{2}h=4a\)，所以 \(AF_1-AF_2=2a\)，故 \(A\) 也在双曲线 \(C\) 的右支上. 因此存在这样的 \(\lambda\)，其值为 \(\frac{12-2\sqrt{2}}{17}\).
\end{enumerate}
\end{solution}
\end{problem}
